\begin{abstract}
Dataset reproducibility depends on documentation completeness, yet no quantitative study has linked repository UX design to metadata outcomes at scale. We introduce friction score (0-4: automated extraction + templates + validation + API) as a quantifiable metric capturing repository design quality.

Through cross-platform analysis of 10,000+ datasets across OpenML, HuggingFace Datasets, and UCI ML Repository, we demonstrate that friction-reduction features correlate with metadata presence. HuggingFace (friction=3) shows 55-66\% optional field presence versus UCI (friction=0) showing 0-15\%, with effect sizes spanning 45-61 percentage points across three fields (preprocessing\_code 61pp, data\_source\_url 51pp, collection\_date 45pp; $p<0.0001$, Cohen's $h$ 1.0-1.8). Within-platform comparison (using synthetic validation data, pending production verification) suggests friction features reduce entry cost: API-uploaded datasets show 12.1pp higher completeness than manual uploads ($p<0.0001$, Cohen's $d=0.621$).

Required fields (license, version) show 74-95\% presence across platforms with weaker friction effects (15-21pp HF-UCI difference, Cramér's $V$ 0.1-0.2), validating enforcement as a distinct but coupled mechanism. Our correlational findings suggest repository design implications: combine enforcement (required fields) with friction-reduction UX (optional fields). Gradient analysis suggests templates may contribute largest marginal effect (estimated 25-30pp) over API (10-15pp) over validation (5-10pp), though feature ablation experiments are needed to confirm these decomposed effects.

This work provides the first quantitative evidence linking friction reduction to metadata completeness at 10,000+ dataset scale, suggesting design hypotheses for controlled deployment testing.
\end{abstract}
